% test-024.tex -- spfuncomp: composicion lineal
%
% Compilar:  lualatex-dev test-024.tex
% Validar:   show-pdf-tags --xml test-024.pdf | rnv-wrapp
%            verapdf --flavour ua2 --format text test-024.pdf
%
% \spfuncomp escribe la composicion lineal f\circ g, (f\circ g) y
% (f\circ g)(x). Intent propio: _compuesta-con:infix(...).
% Las entradas invalidas (que detienen la compilacion) no van aqui.
%
% Los "doc:" son lecturas reales de NVDA con MathCAT (2026-10-08).
\DocumentMetadata
  {
    lang=es-CL,
    pdfversion=2.0,
    pdfstandard=ua-2,
    tagging=on,
    tagging-setup={math/setup={mathml-SE}},
  }
\documentclass{article}
\usepackage[spanish]{babel}
\usepackage{unicode-math}
\usepackage{spintent}
\usepackage{hyperref}
\hypersetup
  {
    pdfdisplaydoctitle = true,
    pdftitle           = {test-024: spfuncomp},
  }
\begin{document}

\section*{A. Composición lineal}

% doc: «f compuesta con g»
1. Dos funciones: $\spfuncomp{f\circ g}$.

% doc: «f compuesta con g compuesta con h»
2. Tres funciones: $\spfuncomp{f\circ g\circ h}$.

% doc: «f compuesta con g compuesta con h compuesta con k»
3. Cuatro funciones: $\spfuncomp{f\circ g\circ h\circ k}$.

% doc: «f compuesta con g»
4. Con espacios: $\spfuncomp{ f \circ g }$.

% doc: «f compuesta con g»
5. Con el carácter Unicode: $\spfuncomp{f ∘ g}$.

% doc: «f sub uno compuesta con g sub dos»
6. Con subíndices: $\spfuncomp{f_{1}\circ g_{2}}$.

% doc: «f prima compuesta con g doble prima»
7. Con primas: $\spfuncomp{f'\circ g''}$.

% doc: «F compuesta con G»
8. Con mayúsculas: $\spfuncomp{F\circ G}$.

% doc: «f inversa compuesta con g»
9. Con la inversa: $\spfuncomp{f^{-1}\circ g}$.

% doc: «f inversa compuesta con g inversa»
10. Dos inversas: $\spfuncomp{f^{-1}\circ g^{-1}}$.

\section*{B. Entre paréntesis}

% doc: «f compuesta con g»
11. Entre paréntesis: $\spfuncomp{(f\circ g)}$.

% doc: «f compuesta con g compuesta con h»
12. Tres entre paréntesis: $\spfuncomp{(f\circ g\circ h)}$.

\section*{C. Aplicada a argumentos}

% doc: «f compuesta con g de x»
13. Aplicada a x: $\spfuncomp{(f\circ g)(x)}$.

% doc: «f compuesta con g de x»
14. Con espacios: $\spfuncomp{( f \circ g ) ( x )}$.

% doc: «f compuesta con g de x coma y»
15. Dos argumentos: $\spfuncomp{(f\circ g)(x,y)}$.

% doc: «f compuesta con g compuesta con h de x»
16. Tres funciones: $\spfuncomp{(f\circ g\circ h)(x)}$.

% doc: «f inversa compuesta con g de x»
17. Con la inversa: $\spfuncomp{(f^{-1}\circ g)(x)}$.

% doc: «f compuesta con g de x más 1»
18. Argumento expresión: $\spfuncomp{(f\circ g)(x+1)}$.

% doc: «f compuesta con g de h de x»
19. Argumento llamada: $\spfuncomp{(f\circ g)(h(x))}$.

% doc: «f compuesta con g de h de k de x»
20. Llamada anidada: $\spfuncomp{(f\circ g)(h(k(x)))}$.

% doc: «f compuesta con g de h inversa de y»
21. Llamada con inversa: $\spfuncomp{(f\circ g)(h^{-1}(y))}$.

% doc: «f compuesta con g de x coma h de y»
22. Letra y llamada: $\spfuncomp{(f\circ g)(x,h(y))}$.

% doc: «f compuesta con g de x coma y coma z»
23. Tres argumentos: $\spfuncomp{(f\circ g)(x,y,z)}$.

% doc: «f compuesta con g de h de x coma k de y»
24. Dos llamadas: $\spfuncomp{(f\circ g)(h(x),k(y))}$.

% doc: «f inversa compuesta con g de x coma y»
25. Inversa y dos argumentos: $\spfuncomp{(f^{-1}\circ g)(x,y)}$.

\section*{D. Contexto}

% doc: «f compuesta con g de x es igual a f de g de x»
26. Con un igual y una llamada: $\spfuncomp{(f\circ g)(x)} = \spfun{f(g(x))}$.

% doc: «f compuesta con g es igual a k»
27. Función igual a otra: $\spfuncomp{f\circ g} = k$.

% doc: «f compuesta con g de x más h de x»
28. En una suma: $\spfuncomp{(f\circ g)(x)} + \spfun{h(x)}$.

% doc: «f compuesta con g de x mayor que 0»
29. Con una desigualdad: $\spfuncomp{(f\circ g)(x)} > 0$.

\section*{E. Llaves}

% doc: silencio
30. silent true: $\spfuncomp[silent=true]{f\circ g}$.

% doc: silencio
31. silent true aplicada: $\spfuncomp[silent=true]{(f\circ g)(x)}$.

% doc: «es igual a 5»
32. silent true y algo después: $\spfuncomp[silent=true]{(f\circ g)(x)} = 5$.

% doc: «f compuesta con g de x»
33. silent false, el valor por defecto: $\spfuncomp[silent=false]{(f\circ g)(x)}$.

\end{document}
